\documentclass[11pt]{article}

% ============================================================
%  AI Applications in Research: Claude Code for Empirical Work
%  A concept summary inspired by Scott Cunningham's "Scott's
%  Mixtape" Substack series on using AI agents for research.
% ============================================================

\usepackage[utf8]{inputenc}
\usepackage[T1]{fontenc}
\usepackage[margin=1in]{geometry}
\usepackage{enumitem}
\usepackage{hyperref}
\usepackage{titlesec}
\usepackage{xcolor}

\hypersetup{
    colorlinks=true,
    linkcolor=blue!50!black,
    urlcolor=blue!60!black,
    citecolor=blue!50!black,
    pdftitle={AI Applications in Research: Claude Code for Empirical Work},
    pdfauthor={Concept Summary}
}

\titleformat{\section}{\large\bfseries}{\thesection}{1em}{}
\setlist[itemize]{leftmargin=1.4em, itemsep=2pt, topsep=3pt}

\title{\textbf{AI Applications in Research}\\[0.3em]
\large Large Language Model Agents and \texttt{Claude Code} in Empirical Social Science}
\author{Concept Summary}
\date{June 30, 2026}

\begin{document}
\maketitle

\begin{abstract}
\noindent This document summarizes the main concepts behind the use of artificial
intelligence---and in particular large language model (LLM) coding agents such as
Anthropic's \texttt{Claude Code}---in academic and applied research. It draws on
the framing developed in Scott Cunningham's \emph{Scott's Mixtape} Substack series
on AI for quantitative social science, together with the econometric framework of
Ludwig, Mullainathan, and Rambachan (2026). The goal is a scannable reference to
how AI agents change the research workflow, the methodological safeguards they
require, and the new opportunities and pitfalls they introduce.
\end{abstract}

% ------------------------------------------------------------
\section{Why AI Now Matters for Research}
% ------------------------------------------------------------

LLMs and agentic coding tools have moved from novelty to infrastructure in
empirical research. Three shifts drive this:

\begin{itemize}
  \item \textbf{Lowered fixed cost of starting.} Projects that were previously too
        costly to attempt---e.g.\ re-classifying hundreds of thousands of historical
        texts---become feasible, so the set of askable questions expands.
  \item \textbf{Agentic execution.} Tools like \texttt{Claude Code} do not merely
        autocomplete code; they read repositories, run scripts, inspect outputs,
        and iterate, acting as a junior collaborator that executes a plan.
  \item \textbf{Measurement at scale.} LLMs can turn unstructured text into
        structured economic variables (sentiment, topic, classification), enabling
        new datasets from previously unusable sources.
\end{itemize}

% ------------------------------------------------------------
\section{What \texttt{Claude Code} Is}
% ------------------------------------------------------------

\texttt{Claude Code} is a command-line AI agent that operates directly inside a
researcher's project directory. In a research setting its distinguishing features are:

\begin{itemize}
  \item \textbf{Repository awareness.} It reads existing files, data, and code to
        ground its actions in the actual project rather than generic boilerplate.
  \item \textbf{Tool use and execution.} It can run shell commands, execute
        statistical scripts (R, Stata, Python), and read back results to verify them.
  \item \textbf{Multi-step orchestration.} It plans, edits across files, and audits
        its own pipeline---including replicating an analysis from scratch in multiple
        languages and packages as a cross-check.
  \item \textbf{Conversational steering.} The researcher iterates with it as a
        thinking partner, refining conjectures rather than issuing one-shot prompts.
\end{itemize}

% ------------------------------------------------------------
\section{The Research Workflow with an AI Agent}
% ------------------------------------------------------------

A recurring theme in Cunningham's series is that the agent is a \emph{collaborator
to be steered}, not an autonomous oracle. The practical workflow:

\begin{enumerate}[leftmargin=1.6em]
  \item \textbf{Frame the conjecture.} The human supplies the economic question and
        identification idea; this is irreducibly human work.
  \item \textbf{Delegate execution.} The agent writes and runs the data-cleaning,
        measurement, and estimation code.
  \item \textbf{Iterate and surface ambiguity.} The agent executes when the next step
        is clear and flags ambiguity when it is not, keeping the human in the loop.
  \item \textbf{Audit and replicate.} Cross-language / cross-package replication and
        coding audits guard against silent errors.
  \item \textbf{Validate against ground truth.} Outputs are checked against a
        human-labeled validation sample before being trusted.
\end{enumerate}

% ------------------------------------------------------------
\section{Econometric Framework: Ludwig, Mullainathan \& Rambachan (2026)}
% ------------------------------------------------------------

The most important methodological reference is \emph{Large Language Models: An
Applied Econometric Framework}. It separates two distinct empirical uses of LLMs,
each with its own validity condition:

\begin{itemize}
  \item \textbf{Prediction problems} (forecasting outcomes from text). Valid
        conclusions require \textbf{no training leakage}: the LLM's training data
        must not overlap with the researcher's outcome sample, or apparent predictive
        power is spurious.
  \item \textbf{Estimation / measurement problems} (using an LLM to label or measure
        an economic concept at scale). Valid downstream inference requires combining
        LLM outputs with a \textbf{small human-validated sample} to correct for
        measurement error and deliver consistent, precise estimates.
\end{itemize}

\noindent The unifying lesson: \textbf{LLM output is a noisy measurement, not ground
truth}, and standard econometric care (validation, bias correction, honest
standard errors) still applies.

% ------------------------------------------------------------
\section{New Problems Created by AI Agents}
% ------------------------------------------------------------

Cunningham emphasizes that solving old problems creates new ones---even when total
time spent is not reduced:

\begin{itemize}
  \item \textbf{Verification burden.} Faster code generation shifts the bottleneck
        to checking that the (plausible-looking) output is actually correct.
  \item \textbf{Automation bias.} Confident, fluent answers can lull researchers into
        skipping scrutiny they would apply to a human RA.
  \item \textbf{Reproducibility and provenance.} Agent-driven, conversational
        workflows can obscure exactly what was run; disciplined version control and
        logged pipelines become essential.
  \item \textbf{Skill atrophy vs.\ leverage.} Delegating routine coding risks
        eroding fluency, even as it frees attention for higher-level design.
\end{itemize}

% ------------------------------------------------------------
\section{Illustrative Applications}
% ------------------------------------------------------------

\begin{itemize}
  \item \textbf{Historical text reclassification.} Re-coding $\sim$305,000 political
        speeches from the late 19th century to the present---a project credited as
        only feasible because of \texttt{Claude Code}.
  \item \textbf{Text-to-data measurement.} Using smaller models (e.g.\ \texttt{gpt-4o-mini})
        for large-scale classification, with the agent orchestrating the pipeline.
  \item \textbf{Live ``on-air'' research.} The \emph{Odd Couple} podcast season with
        economist Caitlin Myers conducts a research project live using \texttt{Claude
        Code}, including visualizations of abortion-clinic access in Texas.
  \item \textbf{Coding audits.} Independent re-implementation of a pipeline across
        multiple languages and packages to detect errors.
\end{itemize}

% ------------------------------------------------------------
\section{Best-Practice Principles}
% ------------------------------------------------------------

\begin{itemize}
  \item Treat the LLM as a \textbf{thinking partner you steer}, not an autonomous agent.
  \item Always pair automated labeling with a \textbf{human validation sample}.
  \item Guard against \textbf{training-data leakage} in any predictive claim.
  \item Maintain \textbf{version control and logged, reproducible pipelines}.
  \item Keep the \textbf{economic question and identification strategy human-owned}.
  \item Build in \textbf{independent replication and audits} rather than trusting first output.
\end{itemize}

% ------------------------------------------------------------
\section*{Selected Resources}
% ------------------------------------------------------------
\addcontentsline{toc}{section}{Selected Resources}

\begin{itemize}
  \item Scott's Mixtape Substack --- \url{https://causalinf.substack.com/}
  \item ``Claude Code'' series index --- \url{https://causalinf.substack.com/s/claude-code}
  \item Ludwig, Mullainathan \& Rambachan, \emph{Large Language Models: An Applied
        Econometric Framework} (NBER WP 33344; arXiv:2412.07031) ---
        \url{https://arxiv.org/abs/2412.07031}
  \item \texttt{MixtapeTools} repository (AI-assisted research templates) ---
        \url{https://github.com/scunning1975/MixtapeTools}
  \item Claude Code documentation --- \url{https://docs.claude.com/en/docs/claude-code}
\end{itemize}

\vfill
\noindent\rule{\linewidth}{0.4pt}\\[2pt]
{\small Prepared as a concept summary of AI applications in research. Newsletter
content is paraphrased from publicly indexed descriptions of Scott Cunningham's
\emph{Scott's Mixtape} Substack; the econometric framework follows Ludwig,
Mullainathan, and Rambachan (2026).}

\end{document}
